\documentclass[a4paper,11pt]{article}
\usepackage{amsmath,amssymb,amsthm}
\usepackage[margin=2.5cm]{geometry}
\usepackage{enumitem}
\usepackage{xcolor}

\newtheorem{proposition}{Proposition}
\newcommand{\open}[1]{\par\smallskip\noindent\fcolorbox{red}{red!5}{\parbox{0.96\linewidth}{\textbf{Open question #1}}}\par\smallskip}
\newcommand{\code}[1]{\texttt{#1}}
\newcommand{\vv}{\mathbf{v}}
\newcommand{\EE}{\mathbf{e}}
\newcommand{\BB}{\mathbf{b}}
\newcommand{\jj}{\mathbf{j}}
\newcommand{\Sm}{\mathsf{S}}
\newcommand{\Pm}{\mathsf{P}}
\newcommand{\Rm}{\mathsf{R}}
\newcommand{\Cm}{\mathsf{C}}
\newcommand{\Gm}{\mathsf{G}}
\newcommand{\Mm}{\mathsf{M}}
\newcommand{\Wb}{\mathsf{W}_B}
\newcommand{\Wt}{\widetilde{\mathsf{W}}}
\newcommand{\Am}{\mathsf{A}}
\newcommand{\Dm}{\mathsf{D}}
\newcommand{\Nset}{\mathcal{N}}
\newcommand{\Nexp}{\widetilde{\mathcal{N}}}
\newcommand{\Hset}{\mathcal{H}}
\newcommand{\Cset}{\mathcal{L}}

\title{An energy-conserving composite-grid ECSIM field solve\\ for AMR in FLEKS}
\author{Derivation note (draft for checking)}
\date{19 November 2026}

\begin{document}
\maketitle

\begin{abstract}
FLEKS solves the implicit ECSIM field equation level by level. Every AMR run of
a uniform plasma at rest loses 0.6--0.9\% of the total energy by $t=1200$, while
the single-level run conserves it to round-off. This note derives one implicit
solve for $\mathbf{E}^{n+\theta}$ on the composite grid of all levels, with the
complete mass matrix. With it the discrete energy identity of single-grid ECSIM
holds across the resolution change, to the tolerance of the linear solver.
Section~\ref{sec:single} restates the single-grid proof in the matrix form used
later. Sections~\ref{sec:current}--\ref{sec:proof} give the composite unknowns,
the prolongation $\Pm$, the composite curl pair and the proof. Sections
\ref{sec:matvec}--\ref{sec:replace} give the matrix-free implementation and list
what it replaces. Section~\ref{sec:div} covers $\nabla\cdot\mathbf{B}$ and Gauss's
law, and Section~\ref{sec:tests} gives the test plan. The open choices are boxed
in red and collected in Section~\ref{sec:open}. The algebraic claims were checked
numerically on a toy composite grid with the FLEKS stencils
(\code{doc/CompositeECSIM\_check.py}). The results are quoted where they are used.
\end{abstract}

\section*{Conventions}
Code units, $c=1$, Gaussian form as in \code{doc/Algorithm.tex}. Levels
$\ell=0,\dots,L$, refinement ratio $r$ (examples use $r=2$), cell volume
$\Delta V_\ell$, spatial dimension $d$ ($d=2$ means $\partial_z\equiv0$ as in
FLEKS 2D). $\mathbf{E}$ lives at nodes and $\mathbf{B}$ at cell centres. The
particle shape is cloud-in-cell (CIC, \code{find\_node\_interpolation}, two nodes
per direction). Vectors of nodal values are bold lower case ($\EE$, $\jj$). Each
node carries 3 components; a nodal scalar operator acts on each component.
$\delta \equiv \theta\Delta t$.

%=====================================================================
\section{Single-grid ECSIM in FLEKS, in discrete form}\label{sec:single}

\subsection{Discrete operators}
Let $a_k$ ($k=x,y,z$) be the node-to-cell difference of FLEKS: at cell $c$ with
corner nodes $c+o$, $o\in\{0,1\}^d$,
\begin{equation}
(a_k f)_c=\frac{1}{2^{d-1}\Delta x_k}\sum_{o,\,o_k=1}\bigl(f_{c+o}-f_{c+o-\hat k}\bigr),
\label{eq:ak}
\end{equation}
i.e.\ the forward difference along $k$ averaged over the $2^{d-1}$ edges of the
cell. This is \code{grad\_node\_to\_center}, \code{div\_node\_to\_center} and
\code{curl\_node\_to\_center} (\code{src/GridUtility.cpp}). For a $Q_1$
(multilinear) interpolant on the cell, $(a_kf)_c$ is the value of $\partial_kf$ at
the cell centre, which equals its average over the cell. The Faraday curl is
\begin{equation}
\Cm=\begin{pmatrix}0&-a_z&a_y\\ a_z&0&-a_x\\ -a_y&a_x&0\end{pmatrix}:
\ \text{nodes}\to\text{cells}.
\end{equation}
The cell-to-node curl \code{curl\_center\_to\_node} uses the $2^d$ cells around
a node with the same $1/2^{d-1}$ average and the backward difference, so
componentwise $\partial^{c\to n}_k=-a_k^{T}$. Inserting this into the components gives
\begin{equation}
\code{curl\_center\_to\_node}=\Cm^{T}\qquad\text{(on a uniform grid)}.
\label{eq:curlT}
\end{equation}
For example, the $y$ component of $\Cm^{T}\BB$ is $-a_x^TB_z+a_z^TB_x
=\partial^{c\to n}_xB_z\cdot(-1)+\partial^{c\to n}_zB_x$, which is the $y$ component
of \code{curl\_center\_to\_node}.

\paragraph{The curl--curl term in the matvec.} \code{update\_E\_matvec} does not
apply $\Cm^T\Cm$. It applies $-\nabla^2+\nabla\nabla\cdot$ as
\code{lap\_node\_to\_node} ($=-\sum_k a_k^Ta_k$) and
\code{grad\_center\_to\_node}$\circ$\code{div\_node\_to\_center}
($=-a_i^T\sum_k a_k$). The two agree when $a_i^Ta_k=a_k^Ta_i$. On a uniform
periodic grid the symbols are
$\hat a_k=e^{i\phi}\,2i\sin(\kappa_k/2)\prod_{k'\ne k}\cos(\kappa_{k'}/2)/\Delta x_k$,
with the common phase $\phi=\sum_k\kappa_k/2$. So $\overline{\hat a_i}\hat a_k$ is real and
symmetric, and
\begin{equation}
-\code{lap}+\code{grad}\,\code{div}=\Cm^T\Cm\quad\text{exactly, on a uniform periodic grid, when \code{coefDiff}}=0.
\end{equation}
With $\code{coefDiff}$ ($=$\code{ratioDivC2C}) $\neq0$, part of $\nabla\cdot$ is
replaced by the centre-to-centre divergence. The operator is then no longer
$\Cm^T\Cm$, and the energy identity below fails even on one level. The AMR
tests already use \code{ratioDivC2C 0}. The composite solve must apply the
curl--curl term as $\Cm^T\Wb\Cm$ (Section~\ref{sec:matvec}). On a uniform grid
this is the same operator, so the single-level results do not change.

\subsection{Particles: gather, deposit, mass matrix}
Particle $p$ (charge $q_p$, mass $m_p$, position $x_p$, the same position for the
mass matrix and for the gather in the mover) has CIC weights $s_{p,g}$ on the
$2^d$ corners $g$ of its cell. Write $\Sm_p$ for the row vector of these weights.
With $\alpha_p=q_p\Delta t/(2m_p)$, the ECSIM velocity update
$\bar\vv_p=\vv_p^n+\alpha_p(\mathbf{E}^{n+\theta}(x_p)+\bar\vv_p\times\mathbf{B}_p)$
is linear in $\mathbf{E}$ for given $\mathbf{B}_p$:
\begin{equation}
\bar\vv_p=\hat\vv_p+\beta_p\,\mathbf{E}^{n+\theta}(x_p),\qquad
\mathbf{E}^{n+\theta}(x_p)=\Sm_p\EE_\theta ,
\label{eq:vbar}
\end{equation}
where $\beta_p$ is the $3\times3$ rotation matrix built in
\code{calc\_mass\_matrix} (\code{alpha[9]}, without the $1/\Delta V$ factor) and
$\hat\vv_p$ is the explicit part (\code{currents[]}$/q_p$). Define the
\emph{unnormalised} deposit and mass matrix
\begin{equation}
\hat\jj=\sum_p q_p\,\Sm_p^T\hat\vv_p ,\qquad
\Mm=\sum_p q_p\,\Sm_p^T\beta_p\Sm_p ,\qquad
\bar\jj=\hat\jj+\Mm\EE_\theta=\sum_pq_p\Sm_p^T\bar\vv_p .
\label{eq:jM}
\end{equation}
FLEKS stores $\code{jHat}=\hat\jj/\Delta V$ and $\code{nodeMM}=\Mm/\Delta V$.
$\Mm$ is not symmetric (the rotation part of $\beta_p$ is antisymmetric), and the proof
does not need it to be.

\paragraph{Kinetic energy.} Since $\vv^{n+1}_p=2\bar\vv_p-\vv^n_p$, we have
$\tfrac12m_p(|\vv^{n+1}_p|^2-|\vv^n_p|^2)=m_p(\vv^{n+1}_p-\vv^n_p)\cdot\bar\vv_p
=2m_p\alpha_p\mathbf{E}(x_p)\cdot\bar\vv_p$, because $(\bar\vv\times\mathbf{B})\cdot\bar\vv=0$.
Summing over the particles and using \eqref{eq:vbar}, \eqref{eq:jM}:
\begin{equation}
K^{n+1}-K^n=\Delta t\sum_pq_p\bar\vv_p\cdot\Sm_p\EE_\theta=\Delta t\,\EE_\theta^T\bar\jj .
\label{eq:dK}
\end{equation}
This is exact for any $\mathbf{B}_p$. The only requirements are that the gather
and the deposit use the same $\Sm_p$ on the same nodal vector $\EE_\theta$, and that
\eqref{eq:vbar} is the velocity the mover actually uses.

\subsection{Fields and the energy identity}
With $\Gm$ the symmetric positive definite field mass matrix ($\Gm=\Delta V\,\mathsf I$
on one level) and $\Wb=\Delta V\,\mathsf I$ on cells, the scheme in weak form is
\begin{align}
\Gm(\EE^{n+1}-\EE^n)&=\Delta t\bigl(\Cm^T\Wb\BB_\theta-4\pi\bar\jj\bigr),\label{eq:amp}\\
\BB^{n+1}-\BB^n&=-\Delta t\,\Cm\EE_\theta ,\label{eq:far}
\end{align}
with $\EE_\theta=\theta\EE^{n+1}+(1-\theta)\EE^n$ and
$\BB_\theta=\BB^n-\theta\Delta t\Cm\EE_\theta$. Divided by $\Delta V$, this is
exactly what \code{update\_E\_rhs}/\code{update\_E\_matvec} solve:
\begin{equation}
\bigl[\mathsf I+\delta^2\Cm^T\Cm+4\pi\delta\,\Mm/\Delta V\bigr]\EE_\theta
=\EE^n+\delta\bigl(\Cm^T\BB^n-4\pi\hat\jj/\Delta V\bigr).
\end{equation}
The energies are $W_E=\EE^T\Gm\EE/8\pi$ and $W_B=\BB^T\Wb\BB/8\pi$
(\code{calc\_E\_field\_energy} and \code{calc\_B\_field\_energy}; the first averages
$|E|^2$ over the corners of each cell, which gives $\Gm=\Delta V\,\mathsf I$).

\begin{proposition}\label{prop:energy}
Let $\Gm$ be symmetric and let the Ampère curl be $\Cm^T\Wb$ with the same $\Cm$
as in Faraday's law \eqref{eq:far}. Then \eqref{eq:dK}--\eqref{eq:far} give
\begin{equation}
(W_E+W_B+K)^{n+1}-(W_E+W_B+K)^n=-\frac{\theta-\frac12}{4\pi}
\Bigl[(\Delta\EE)^T\Gm\Delta\EE+(\Delta\BB)^T\Wb\Delta\BB\Bigr],
\label{eq:energy}
\end{equation}
where $\Delta\EE=\EE^{n+1}-\EE^n$ and $\Delta\BB=\BB^{n+1}-\BB^n$.
For $\theta=\frac12$ the total energy is conserved exactly.
\end{proposition}
\begin{proof}
$\EE^{n+1}+\EE^n=2\EE_\theta-(2\theta-1)\Delta\EE$. With $\Gm$ symmetric,
$8\pi\Delta W_E=\Delta\EE^T\Gm(\EE^{n+1}+\EE^n)
=2\EE_\theta^T\Gm\Delta\EE-(2\theta-1)\Delta\EE^T\Gm\Delta\EE$. By \eqref{eq:amp},
$\EE_\theta^T\Gm\Delta\EE=\Delta t\,(\Cm\EE_\theta)^T\Wb\BB_\theta-4\pi\Delta t\,\EE_\theta^T\bar\jj$.
In the same way, $\BB^{n+1}+\BB^n=2\BB_\theta-(2\theta-1)\Delta\BB$ and \eqref{eq:far}
give $8\pi\Delta W_B=-2\Delta t\,(\Cm\EE_\theta)^T\Wb\BB_\theta-(2\theta-1)\Delta\BB^T\Wb\Delta\BB$.
Adding these, the two curl terms cancel identically. The current term,
$-\Delta t\,\EE_\theta^T\bar\jj$, cancels \eqref{eq:dK}.
\end{proof}

\paragraph{Effect of the solver tolerance.} If the linear system is solved with
residual $\mathbf r$ (left side minus right side of the $\Gm$-weighted system
\eqref{eq:Asys} below, $\mathbf r=\Am\EE_\theta-\text{rhs}$), then \eqref{eq:amp} holds with an extra $+\mathbf r/\theta$ on the
right. The total energy changes per step by the extra amount
\begin{equation}
\delta W=\frac{1}{4\pi\theta}\,\EE_\theta^T\mathbf r .
\label{eq:tol}
\end{equation}
So ``exact to the solver tolerance'' means $|\delta W|\le|\EE_\theta|\,|\mathbf r|/(4\pi\theta)$
in the $\Gm$-consistent norms (Section~\ref{sec:precond}).

\paragraph{Conditions used} (each must survive the move to AMR):
\begin{enumerate}[label=(C\arabic*)]
\item gather and deposit use the same weights on the same unknowns;
\item $\bar\jj=\hat\jj+\Mm\EE_\theta$ is the current of the velocities the mover
  produces (no smoothing of $\hat\jj$, $\Mm$ or $\mathbf{E}$ unless it is applied as
  $\Sm\to\Sm\mathsf F$ consistently in all three places);
\item the Ampère curl is the $(\Gm,\Wb)$-adjoint of the Faraday curl;
\item $\Gm$ is symmetric, and $\EE^T\Gm\EE/8\pi$ is the energy that is reported;
\item no extra terms in the operator: upwind $\mathbf{E}$ off, $\code{coefDiff}=0$,
  div~$\mathbf{E}$ correction off (it moves the particles), $\theta=\frac12$ for
  exact conservation.
\end{enumerate}

%=====================================================================
\section{The current AMR algorithm and where the identity fails}\label{sec:current}

This section describes the full-PIC ECSIM path (\code{Pic::update} with
\code{solveEM}) as implemented now.
\begin{itemize}
\item \textbf{Deposit} (\code{calc\_mass\_matrix\_amr},
  \code{Particles::calc\_mass\_matrix\_amr}). A particle on level $\ell$ deposits
  $\hat\jj$ and $\Mm$ with the level-$\ell$ CIC shape on level $\ell$. It also deposits
  with the coarse shapes on \emph{every coarser level} (\code{jhc}, \code{nmmc}).
  If its tile has the \code{refined\_neighbour} flag, it deposits with the
  level-$(\ell+1)$ shape on level $\ell+1$ as well (\code{jhf}, \code{nmmf}). Each
  level therefore sees all the plasma inside its own region, covered part included.
\item \textbf{Solve} (\code{update\_E\_impl}). The levels are solved one after the
  other, coarse first. The unknowns of level $\ell$ are its owned valid nodes. On
  a fine level these include the edge nodes on the coarse--fine interface, so the
  interface carries two independent values of $\mathbf{E}$, one per level.
  The fine ghost nodes are a Dirichlet condition: in the initial residual
  (\code{useZeroBC=false}) they are the bilinear interpolation of the coarse
  $\mathbf{E}^{n+\theta}$ (\code{fill\_fine\_lev\_bny\_from\_coarse} on
  \code{nodeEth[iLev-1]}), and inside GMRES they are zero.
  The covered coarse nodes are unknowns of the coarse solve.
\item \textbf{Ampère curl} (\code{update\_E\_rhs}). This is \code{curl\_center\_to\_node}
  of \code{centerB}, with the fine ghost cells interpolated from the coarse $\mathbf{B}$.
\item \textbf{Faraday} (\code{update\_B}). Each level applies \code{curl\_node\_to\_center} to its
  own \code{nodeEth}, the covered coarse cells included. Covered $\mathbf{B}$ is
  averaged down only if \code{projectDownEmFields} is set.
\item \textbf{Mover}. A particle on level $\ell$ gathers \code{nodeEth[$\ell$]} with its CIC
  weights on the corners of its own (leaf) cell.
\item \textbf{Diagnostic}. $W_E$ and $W_B$ are summed over the unrefined cells of every
  level, with each level's own nodal $\mathbf{E}$.
\end{itemize}
Measured against conditions (C1)--(C4):
\begin{description}
\item[(F1)] The cross-level deposits violate (C1) and (C2). A fine particle puts
  $\Mm$ on the coarse level, but it never feels the coarse $\mathbf{E}$. A coarse particle
  next to the interface puts $\Mm$ on the fine level, but it gathers only from the
  coarse level. The coarse $\Mm$ is therefore not the response of any particle to
  the coarse $\mathbf{E}$ that the coarse solve determines. This is the diagnosis
  that the mass matrix does not act across the resolution change.
\item[(F2)] The two values of $\mathbf{E}$ on the interface violate (C1) and (C3):
  the coarse leaf cells and the fine cells next to the interface are advanced with
  different values of $\mathbf{E}$.
\item[(F3)] The Ampère curl at the fine edge nodes uses ghost $\mathbf{B}$
  interpolated from the coarse level. It is not the transpose of the Faraday curl
  that advances the leaf cells, which violates (C3).
\item[(F4)] The covered coarse $\mathbf{E}$ and $\mathbf{B}$ are evolved. They are
  not counted in the energy, but they feed the coarse interface nodes.
\end{description}
The hybrid-only routines \code{sync\_emf\_at\_interface}, \code{fill\_ghost\_emf\_from\_coarse},
the Faraday-evolved ghost layer (\code{faraday\_ngrow}) and
\code{relax\_covered\_B\_to\_fine} are in \code{PicHybrid.cpp}. The ECSIM path does not
call them. Section~\ref{sec:replace} says what the composite operators make of them.

%=====================================================================
\section{Composite unknowns}\label{sec:unknowns}

\paragraph{Leaf cells.} $\Cset$ is the set of cells not covered by a finer level,
with $\Wb=\mathrm{diag}(\Delta V_c)$, $c\in\Cset$. This is the set over which
\code{calc\_B\_field\_energy} sums. Each particle lives in exactly one leaf cell,
because AMReX puts a particle on the finest level that contains it.

\paragraph{Expanded node set $\Nexp$.} This is the set of corners of the leaf cells.
Because the shape is CIC, \emph{every particle stencil and every leaf-cell curl lies
in $\Nexp$, and no ghost node of any level enters.} $\Nexp$ consists of
\begin{itemize}
\item the nodes of each level that are interior to its leaf region;
\item the \emph{coincident} interface nodes: fine nodes on the coarse--fine
  interface that coincide with a coarse node. They are corners of both coarse and
  fine leaf cells;
\item the \emph{hanging} nodes $\Hset$: fine nodes on the interface that do not
  coincide with a coarse node (midpoints of coarse edges, centres of coarse faces).
  They are corners of fine leaf cells only.
\end{itemize}
The \emph{covered coarse nodes}, which lie strictly inside the region covered by
level $\ell+1$, are not corners of any leaf cell. They are not in $\Nexp$.

\paragraph{Unknowns.} $\Nset=\Nexp\setminus\Hset$, one value per physical location.
The hanging nodes are slaved by $\Pm$ (Section~\ref{sec:P}). The covered coarse nodes
are \emph{eliminated}: they appear in none of the equations. After the solve they are
filled by a restriction $\Rm$ (from the fine $\mathbf{E}$ by injection or full weighting)
for output, coupling, the coarse div~$\mathbf{E}$ diagnostics and regridding. They
enter neither the energy nor the dynamics. (The alternative of keeping them as
unknowns tied by $\EE_c=\Rm\EE_f$ adds constraint rows and changes nothing physical,
because no leaf cell or particle uses them.)

\open{Q1 (storage of the coincident interface unknowns).} Either the finer level owns
the coincident nodes (recommended) or the coarser level does. If the finer level owns
them, the hanging values are averages of neighbouring fine edge nodes, so $\Pm$ is
local to the fine level (Section~\ref{sec:P}), and the coarse copies are an injection,
which is the adjoint of a plain sum. If the coarser level owns them, $\Pm$ becomes a
coarse-to-fine interpolation like the present ghost fill, but every matvec then needs
two-way coarse--fine communication.

\open{Q2 (restriction $\Rm$ for the covered coarse $\mathbf{E}$).} Injection or full
weighting. It is a diagnostic choice only. Full weighting is smoother for
output and coupling to MHD.

%=====================================================================
\section{Prolongation, deposit, mass matrix and field mass matrix}\label{sec:P}

\paragraph{Prolongation.} $\Pm:\mathbb{R}^{3|\Nset|}\to\mathbb{R}^{3|\Nexp|}$ is the
identity on $\Nset$. At a hanging node it is the multilinear interpolation, along the
interface, from the coincident nodes of the coarse edge or face that contains it:
weight $1/2$ from each end of an edge (edge midpoint) and $1/4$ from each corner of a face
(face centre), for $r=2$. For general $r$ the weights are $j/r$ products. $\Pm$ makes
$\Pm\EE$ the restriction of a continuous piecewise-$Q_1$ field. On a coarse--fine face
the fine piecewise-bilinear trace then equals the coarse bilinear trace exactly.

\paragraph{Gather and deposit.} For particle $p$ in leaf cell $c(p)$, $\Sm_p$ are its CIC
weights on the corners of $c(p)$, as row vectors on $\Nexp$. Then
\begin{equation}
\mathbf{E}^{n+\theta}(x_p)=\Sm_p\Pm\EE_\theta ,\qquad
\hat\jj=\Pm^T\sum_pq_p\Sm_p^T\hat\vv_p=\Pm^T\hat{\tilde\jj} ,\qquad
\Mm=\Pm^T\Bigl(\sum_pq_p\Sm_p^T\beta_p\Sm_p\Bigr)\Pm=\Pm^T\widetilde\Mm\,\Pm .
\label{eq:PMP}
\end{equation}
Here $\hat{\tilde\jj}$ and $\widetilde\Mm$ are the per-level deposits of the particles
on their own level only, on the corners of their own cells, hanging nodes
included. These are the \code{calc\_mass\_matrix} deposits per level
\emph{without} the cross-level parts \code{jhc}, \code{nmmc}, \code{jhf} and \code{nmmf}.
$\Pm^T$ folds the hanging rows and columns onto the coincident nodes with the
interpolation weights. The only particles affected are those in the first fine
cell layer next to the interface, whose cells have hanging corners. Coarse particles
next to the interface use only the coarse corners, which are coincident
nodes. With \eqref{eq:PMP}, \eqref{eq:dK} holds unchanged:
$\Delta K=\Delta t\,\EE_\theta^T\bar\jj$.

\paragraph{Field mass matrix.} Let $\Wt=\mathrm{diag}(\tilde w_g)$ on $\Nexp$, with
$\tilde w_g=\sum_{c\in\Cset,\,g\in c}\Delta V_c/2^d$. Then
\begin{equation}
\Gm=\Pm^T\Wt\Pm .
\label{eq:G}
\end{equation}
$\EE^T\Gm\EE/8\pi$ is \emph{exactly} what \code{calc\_E\_field\_energy} computes,
provided the level arrays hold the composite field: hanging nodes $=\Pm\EE$, and the
coincident nodes equal on both levels. $\Gm$ is symmetric positive definite. It is
diagonal except for couplings among the coincident nodes of one coarse face, which
come through the hanging nodes.

\open{Q3 (consistent or lumped $\Gm$).} With the lumped $\Gm_L=\mathrm{diag}(\Gm\mathbf 1)$,
Proposition~\ref{prop:energy} still holds, but for the energy $\EE^T\Gm_L\EE/8\pi$,
which is not what the diagnostic reports. In the toy check (one step, random
fields), the lumped scheme conserves its own energy to $2\times10^{-16}$ but changes the
diagnostic energy by $5.6\times10^{-3}$ (2D) and $2.3\times10^{-3}$ (3D) relative. The
consistent $\Gm$ conserves the diagnostic energy to $10^{-16}$. Recommendation:
use the consistent $\Gm$. It costs one extra $\Pm$/$\Pm^T$ pair in the matvec, which is
needed for $\Mm$ anyway.

%=====================================================================
\section{Composite curl pair and the energy proof}\label{sec:proof}

\paragraph{Faraday.} For every leaf cell $c$ at level $\ell$,
$(\Cm\EE)_c=(\Cm_\ell\,\Pm\EE)_c$, the level-$\ell$ stencil \eqref{eq:ak} on the corners
of $c$:
\begin{equation}
\BB_c^{n+1}=\BB_c^n-\Delta t\,(\Cm_\ell\Pm\EE_\theta)_c,\qquad c\in\Cset .
\end{equation}
Only leaf cells are advanced. Covered coarse cells are restricted for output
(\code{average\_down}). Fine ghost cells are not advanced at all.

\paragraph{Ampère.} The curl of $\mathbf{B}$ in \eqref{eq:amp} is $\Cm^T\Wb\BB$ with
$\Cm=\Cm_{\rm leaf}\Pm$:
\begin{equation}
(\Cm^T\Wb\BB)_g=\sum_\ell\Delta V_\ell\,\bigl[\Pm^T\Cm_\ell^{T}(\chi_\ell\BB_\ell)\bigr]_g ,
\label{eq:ampcurl}
\end{equation}
where $\chi_\ell$ masks the leaf cells of level $\ell$ (zero on covered cells
\emph{and on the ghost cells of level $\ell$}). By \eqref{eq:curlT},
$\Cm_\ell^T(\chi_\ell\BB_\ell)$ is \code{curl\_center\_to\_node} applied to the masked
$\mathbf{B}$. At a node inside one level, \eqref{eq:ampcurl} divided by
$\tilde w_g=\Delta V_\ell$ is the present curl. At a coincident node, the coarse leaf
cells around it enter with weight $\Delta V_0$ and the fine cells with $\Delta V_1$, and
the fine cells around the neighbouring hanging nodes enter through $\Pm^T$. No ghost
$\mathbf{B}$ is needed.

\paragraph{The composite linear system.} Substituting $\bar\jj=\hat\jj+\Mm\EE_\theta$ and
$\BB_\theta=\BB^n-\delta\Cm\EE_\theta$ into \eqref{eq:amp} gives
\begin{equation}
\Am\EE_\theta\equiv\bigl[\Gm+\delta^2\Cm^T\Wb\Cm+4\pi\delta\,\Pm^T\widetilde\Mm\Pm\bigr]\EE_\theta
=\Gm\EE^n+\delta\bigl(\Cm^T\Wb\BB^n-4\pi\Pm^T\hat{\tilde\jj}\bigr),
\label{eq:Asys}
\end{equation}
followed by $\EE^{n+1}=(\EE_\theta-(1-\theta)\EE^n)/\theta$ on $\Nset$, the
hanging nodes by $\Pm$, and Faraday on $\Cset$. The update is linear, so if $\EE^n$
is a composite field (hanging $=\Pm$, coincident nodes single-valued), $\EE^{n+1}$ is one too.

\begin{proposition}\label{prop:comp}
The composite scheme \eqref{eq:PMP}--\eqref{eq:Asys} with Faraday on $\Cset$ satisfies
the energy identity \eqref{eq:energy} with $\Gm=\Pm^T\Wt\Pm$ and $\Wb$ the leaf
volumes. For $\theta=\frac12$ it conserves $W_E+W_B+K$, as computed by the existing
diagnostics, up to the residual term \eqref{eq:tol}.
\end{proposition}
\begin{proof}
The proof of Proposition~\ref{prop:energy} used only these facts: $\Gm$ is symmetric;
the Ampère curl is $\Cm^T\Wb$ for the Faraday $\Cm$; $\Wb$ is symmetric; and
$\Delta K=\Delta t\EE_\theta^T\bar\jj$ with $\bar\jj$ the current in Ampère's law. The
first three hold by construction. The last follows from \eqref{eq:PMP}: the gather
$\Sm_p\Pm$ and the deposit $(\Sm_p\Pm)^T$ are transposes. Nothing in the proof refers
to a level.
\end{proof}

\paragraph{Numerical check.} On a toy periodic grid with one refined box
(\code{CompositeECSIM\_check.py}, $r=2$, CIC, the FLEKS stencils, random $\mathbf{E}^n$
and $\mathbf{B}^n$, 40 particles per leaf cell with random $\mathbf{B}_p$), one $\theta=\frac12$
step with a direct solve changes the total energy by $0$ (2D) and
$1.3\times10^{-16}$ (3D) relative.

%=====================================================================
\section{Matrix-free matvec, right-hand side and preconditioning}\label{sec:matvec}

\subsection{Matvec $\mathbf y=\Am\mathbf x$}
\begin{enumerate}
\item \textbf{Scatter} $\mathbf x\to$ level arrays (\code{solverVecMF[$\ell$]}). Fill the
  owned unknowns, then the coarse copies of the coincident nodes (injection from the
  fine owner, if Q1 is answered ``fine owns''), then the hanging nodes by $\Pm$ (an
  average of fine edge neighbours; fine-level halo exchange only). Set the covered
  coarse nodes and the ghost nodes to \emph{zero}. Their values must not matter, and
  zero makes any leak visible in testing.
\item \textbf{Per level} $\ell$, on level arrays:
  \begin{enumerate}
  \item $\mathbf b_\ell=\chi_\ell\,\code{curl\_node\_to\_center}(\mathbf x_\ell)$;
  \item $\mathbf y_\ell=\Delta V_\ell\bigl[\delta^2\,\code{curl\_center\_to\_node}(\mathbf b_\ell)
    +4\pi\delta\,\code{nodeMM}_\ell\,\mathbf x_\ell+\omega_\ell\,\mathbf x_\ell\bigr]$, where
    $\code{nodeMM}_\ell$ is the own-level deposit (the existing
    \code{update\_E\_M\_dot\_E}) and $\omega_\ell(g)=n_{\rm leaf,\ell}(g)/2^d$ is the
    fraction of the $2^d$ cells around $g$ that are leaf cells of level $\ell$.
    (Then $\sum_\ell\Delta V_\ell\omega_\ell=\tilde w_g$.)
  \end{enumerate}
\item \textbf{Fold} $\Pm^T$. Add each hanging-node $\mathbf y$ to its coincident neighbours
  with the interpolation weights; add the coarse copies of the coincident nodes into
  the fine owner (the adjoint of the injection); discard the covered and ghost values.
  This uses the same \code{ParallelAdd} pattern as the present
  \code{calc\_mass\_matrix\_amr}.
\item \textbf{Row scaling} for GMRES: divide row $g$ by the lumped weight $(\Gm\mathbf 1)_g$.
  In the interior of a level this reproduces the present operator
  $\mathsf I+\delta^2\Cm^T\Cm+4\pi\delta\,\code{nodeMM}$ exactly.
\end{enumerate}
The curl--curl term must be the $\Cm^T(\chi\,\Cm\,\cdot)$ product of steps 2(a) and 2(b),
\emph{not} \code{lap\_node\_to\_node}$+$\code{grad\,div}. The two agree in the interior
of a level, but only the masked curl pair gives the transpose at the interface and at
covered cells. The extra cost over the present level matvecs is one fold and
one scatter per matvec. The total number of unknowns is smaller (no covered coarse
nodes, no duplicated interface).

\paragraph{Right-hand side.} By the same steps:
$\Gm\EE^n+\delta(\Cm^T\Wb\BB^n-4\pi\Pm^T\hat{\tilde\jj})$, with the masked
\code{curl\_center\_to\_node} of \code{centerB} and \emph{no} ghost fill of \code{centerB}
or \code{nodeB}. The co-moving term becomes $4\pi\delta\,\Pm^T\widetilde\Mm\Pm\,\EE_0$. Its effect
on the energy is not analysed here: it is not used in the rest tests.

\paragraph{Physical boundaries.} The present code applies the field BCs (Dirichlet rows,
conducting walls, \code{apply\_field\_bc\_homogeneous}) only on level~0. In the
composite system they apply to the boundary nodes of whichever level owns them.
\open{Q4 (BCs on refined levels touching the domain boundary).} Either require that
refined levels do not touch non-periodic boundaries, or extend the BC rows to all
levels. The tests are periodic, so this can wait.

\subsection{GMRES norm, tolerance and preconditioning}\label{sec:precond}
\begin{itemize}
\item \textbf{Norm.} By \eqref{eq:tol}, the energy error is $\EE_\theta^T\mathbf r$. With the
  row scaling of step 4, the natural inner product is
  $\langle\mathbf u,\mathbf v\rangle=\mathbf u^T\mathrm{diag}(\Gm\mathbf 1)\mathbf v$, i.e.\
  volume weighted. The present unweighted $\ell_2$ norm weights a coarse node like a fine one,
  although the coarse node carries $2^d$ times the volume; the mismatch grows with every level.
\item \textbf{Spectrum.} At fixed $\Delta t$ the fine level has $c\Delta t/\Delta x$ that is
  $r$ times larger. The iteration count should therefore be close to that of a
  single-level run at the fine resolution, not that of the present coarse solve. The
  coarse--fine coupling adds no stiffness of its own, but it does couple the
  long-wavelength coarse modes to the fine solve.
\item \textbf{Preconditioners}, in order of effort: (P0) none, as now, with the row
  scaling, to measure the baseline; (P1) point-block Jacobi with the $3\times3$ node
  blocks of $\Am$, which are available from \code{nodeMM} and the stencil diagonal;
  (P2) level-block additive Schwarz, with a few GMRES iterations of the existing
  level operator as the block solve and flexible GMRES outside. This is not the
  rejected iteration of level solves: the outer composite GMRES makes the result exact,
  and the inner solves are cheap and inexact; (P3) composite multigrid, later.
\end{itemize}
\open{Q5 (preconditioner and tolerance).} Recommendation: start with P0 plus row scaling
and the weighted norm. Measure the energy error against \code{EFieldTol} (it should
scale linearly) and the iteration counts against the present solver, then decide on
P1 or P2.

%=====================================================================
\section{What the composite solve replaces}\label{sec:replace}
In the full-PIC (ECSIM) path:
\begin{enumerate}
\item The level loop in \code{update\_E\_impl} becomes one solve of \eqref{eq:Asys}.
  \code{convert\_1d\_to\_3d} and \code{convert\_3d\_to\_1d} run over all levels, with
  ownership by $\Nset$.
\item The fine Dirichlet ghost $\mathbf{E}$ in \code{update\_E\_matvec}
  (\code{fill\_fine\_lev\_bny\_from\_coarse} of \code{nodeEth[iLev-1]}) is replaced by $\Pm$
  on the hanging nodes. The ghost nodes are no longer used.
\item The ghost fill of \code{centerB} and \code{nodeB} for the Ampère curl in
  \code{update\_E\_rhs} is replaced by the masked transpose \eqref{eq:ampcurl}.
\item The cross-level deposits \code{jhc}, \code{nmmc}, \code{jhf} and \code{nmmf} in
  \code{calc\_mass\_matrix\_amr} are replaced by own-level deposits folded with $\Pm^T$.
  (The \code{sum\_two\_lev\_interface\_node} loops in \code{calc\_mass\_matrix} run only
  when there is a single level, so they do nothing.)
\item The coarse solve and the coarse Faraday update on the covered region are replaced
  by $\Rm$ and \code{average\_down}, for output only. \code{project\_down\_E} becomes $\Rm$.
\item \code{update\_B} advances only the leaf cells. Fine ghost \code{centerB} is
  still needed for \code{average\_center\_to\_node}, because \code{nodeB} at the fine edge nodes
  enters $\beta_p$ and the pusher. This does not affect the energy, since the
  $\mathbf{B}_p$ used is arbitrary in \eqref{eq:dK}, but it must be the same $\mathbf{B}_p$
  in the mass matrix and in the mover.
\end{enumerate}
\open{Q6 (nodal $\mathbf{B}$ at interface nodes).} Either use the volume-weighted average
over the leaf cells around the node (composite, single-valued) or keep the present
fine ghost interpolation. Both conserve energy. The composite average makes the
coarse and fine particles on either side see the same $\mathbf{B}$ at the coincident
nodes.

\paragraph{The hybrid interface machinery.} \code{sync\_emf\_at\_interface},
\code{fill\_ghost\_emf\_from\_coarse}, the Faraday-evolved first ghost layer and
\code{relax\_covered\_B\_to\_fine} belong to the hybrid solver and are not touched. The
composite construction is their ECSIM counterpart. The coincident nodes carry a single
$\mathbf{E}$, which is what the synchronization of $\mathbf{E}$ provides. The hanging nodes
interpolated along the coarse edges are $\Pm$, which is what the synchronization does
on the fine side. The fine cells next to the interface are leaf cells and are advanced
by Faraday's law with that $\mathbf{E}$, so no ghost layer is evolved. The covered cells
are not advanced at all, so they cannot drift, and their restriction is diagnostic only.
If the hybrid solver later adopts $\Pm$ and the leaf-only Faraday update, these four
routines reduce to $\Pm$, nothing, nothing and \code{average\_down}.

%=====================================================================
\section{Divergence and charge}\label{sec:div}

\subsection{$\nabla\cdot\mathbf{B}$}
On a uniform grid the conserved divergence is the weak one, $\Gamma^T\Wb\BB$, where
$\Gamma=(a_x,a_y,a_z)^T$ is the node-to-cell gradient; up to sign and volume this is
\code{div\_center\_to\_node}. It is conserved because $\Gamma^T\Cm=0$: the symbol is
$\bar{\hat a}\cdot(\hat a\times\cdot)$ with $\hat a=e^{i\phi}i\times$(real vector), so
$\bar{\hat a}\times\hat a=0$. On the composite grid the candidate is $\Gamma_P=\Gamma_{\rm leaf}\Pm$
(test functions continuous piecewise $Q_1$), and the question is whether $\Gamma_P^T\Wb\Cm=0$.

\paragraph{2D: exact.} In 2D only $B_x=a_yE_z$ and $B_y=-a_xE_z$ enter. Per cell and unit length in $z$,
$(\Delta V_c/\Delta z)(a_x\phi\,a_y\psi-a_y\phi\,a_x\psi)=\tfrac12[(\phi_{10}-\phi_{01})(\psi_{11}-\psi_{00})
-(\phi_{11}-\phi_{00})(\psi_{10}-\psi_{01})]$. This is the shoelace form of the edge-trapezoid
integral $\oint_{\partial c}\phi\,d\psi$, which holds as an identity in the corner values.
Interior edges cancel. On a coarse edge split by a hanging node with the
$\Pm$ values $\phi_m=\frac{\phi_a+\phi_b}2$, $\psi_m=\frac{\psi_a+\psi_b}2$, the two fine
trapezoids sum to $\frac12(\phi_a+\phi_b)(\psi_b-\psi_a)$, which is the coarse one (the trapezoid
rule is exact for $\int\phi\,d\psi$ with linear $\phi,\psi$; this holds for any $r$). Hence
$\Gamma_P^T\Wb\Cm=0$ exactly in 2D. Numerically: $0$.

\paragraph{3D: not exact.} In 3D the centre value of $a_x\phi$ averages the two $z$-layers
of the cell, so the cell term is the 2D Jacobian of the $z$-averaged fields. Across a
coarse face, the coarse cell has one $z$-midline and the two fine layers have two. A hand count with
$\Pm$ shows that the bilinear edge terms then come with weights $(\frac{10}{16},\frac{6}{16})$ on the fine side
and $(\frac{8}{16},\frac{8}{16})$ on the coarse side, and they do not cancel. In the toy
3D check, $\max|\Gamma_P^T\Wb\Cm|=0.125$, nonzero at exactly the coarse interface nodes
(26 of 26) and zero elsewhere. The left null space of $\Cm$ has dimension 311, against 241
independent nodal test functions. So conserved linear functionals of $\mathbf{B}$ do exist, but
whether a \emph{local} composite divergence among them exists is not known.
\open{Q7 (3D $\nabla\cdot\mathbf{B}$ at the interface).} (a) Accept an interface source
of $\nabla\cdot\mathbf{B}$ and rely on the existing cleaning (\code{correct\_B}, hyperbolic
cleaning). These change $W_B$, so they must be off in the energy tests. (b) Look for a
local functional in the left null space of $\Cm$, i.e.\ a modified weak divergence at the
interface nodes. (c) Modify $\Cm$ at the interface so that $\Gamma_P^T\Wb\Cm=0$. Any
modified $\Cm$ keeps the energy proof, as long as Ampère uses its transpose. The present
scheme does not conserve $\nabla\cdot\mathbf{B}$ at the interface either, so (a) is no
regression.

\subsection{Gauss's law and charge}
ECSIM is not charge conserving, which is why FLEKS has the div~$\mathbf{E}$ correction.
The curl term is clean on one level: \code{div\_node\_to\_center}$\circ\Cm^T=0$ by the same
symbol argument. So $\nabla\cdot\mathbf{E}$ changes only through $\nabla\cdot\bar{\mathbf J}$.
On the composite grid the curl term of Ampère's law is $\Gm^{-1}\Cm^T\Wb\BB$, and
$\Dm_{\rm leaf}\Pm\,\Gm^{-1}\Cm^T\Wb\ne0$ at the interface. In the toy check the maximum
is $0.45$ (2D) and $0.23$ (3D), with the consistent and the lumped $\Gm$ alike. The
residual sits in the cell layers on both sides of the interface (3D toy: all 56 fine and all
56 coarse cells adjacent to it) and is zero elsewhere. So the interface is a source of spurious charge proportional to
the curl of $\mathbf{B}$, even in vacuum.
\open{Q8 (Gauss's law).} The div~$\mathbf{E}$ correction (\code{amr\_divE\_correction}) must
use the composite divergence of $\Pm\mathbf{E}$ on leaf cells and the composite charge
(leaf-cell deposit with the same CIC weights). The particle position correction breaks
the energy identity, because it moves particles after $\Mm$ is formed, so it is off in the
energy tests. Whether a Gauss-compatible $\Gm$ or $\Pm$ exists is open, and it is tied to Q7.

%=====================================================================
\section{Multiple levels and regridding}
With proper nesting, each hanging node of level $\ell+1$ lies on a face of level $\ell$
cells. If that face is itself on the $\ell/\ell-1$ interface, the $\ell$ values used
are composite values, and $\Pm=\Pm_1\Pm_2\cdots$ composes. Nothing in the proof changes.
\open{Q9 (nesting).} Does the FLEKS refinement guarantee a buffer of at least one level-$\ell$
cell between the $\ell+1$ and $\ell-1$ boundaries? If it does not, the composed $\Pm$ at the
places where they touch needs care.

At a regrid the discrete space changes. $\EE$ must be projected onto the new composite
space (hanging $=\Pm$, coincident nodes single-valued), which changes $W_E$ once per regrid.
The energy tests use a static grid.

%=====================================================================
\section{Test plan}\label{sec:tests}
All the tests use $\theta=0.5$, \code{ratioDivC2C 0}, \code{\#DIVE F}, no smoothing of
$\mathbf{E}$ or $\mathbf{J}$, no upwind $\mathbf{E}$, \code{projectDownEmFields} off, a static
grid, and energy from the log every step (\code{\#SAVELOG 1}). Every batch gets $z=0$ maps,
not only tables.
\begin{enumerate}[label=T\arabic*]
\item \textbf{Algebra} (done): \code{CompositeECSIM\_check.py 2|3}. Energy to $10^{-16}$,
  2D weak $\nabla\cdot\mathbf{B}$ exact, the 3D interface $\nabla\cdot\mathbf{B}$ and the
  Gauss residuals as quoted. Rerun it after any change of $\Pm$ or $\Cm$.
\item \textbf{Operator unit test} in FLEKS, before any physics. On a two-level grid,
  compare the matrix-free $\Am$ against an assembled $\Am$ column by column on random
  vectors. Check $\mathbf u^T(\Cm^T\Wb\Cm)\mathbf v=(\Cm\mathbf u)^T\Wb(\Cm\mathbf v)$, check that
  $\Pm^T$ is the adjoint of $\Pm$, and check that $\EE^T\Gm\EE/8\pi$ equals
  \code{calc\_E\_field\_energy}. Also check that the matvec on a single level is identical to
  the present one when $\code{coefDiff}=0$.
\item \textbf{Rest}: a uniform plasma at rest on the AMR grid with equal particle
  weights on the two levels (\code{\#CONSTANTPPV}), to $t=1200$.
  The expected energy error per step is \eqref{eq:tol}. Run \code{EFieldTol} $=10^{-6}$,
  $10^{-8}$ and $10^{-10}$: the drift must scale linearly with the tolerance and must not
  depend on the refinement. Compare with the uniform grid (exact) and with the
  level-by-level solve, split into electrons and ions.
\item \textbf{2D check}: a fake-2D rest run. The weak $\nabla\cdot\mathbf{B}$ at the
  interface nodes should be at round-off. In 3D it should not be, and its size should be
  reported.
\item \textbf{Waves}: fast waves and drift waves crossing the refined box, AMR
  against the uniform coarse and fine grids. Compare the phase and amplitude errors in
  the box and the reflection at the interface with the present solver, and give maps.
  Report the GMRES iterations per step against the present solver.
\item \textbf{$\theta>\frac12$}, optional: the energy must decrease by exactly the right side
  of \eqref{eq:energy}. (Keep in mind the known $P_e$ problems at $\theta=0.51$.)
\end{enumerate}

%=====================================================================
\section{Implementation in FLEKS (23 November 2026)}\label{sec:impl}
The solve is in \code{src/PicComposite.cpp} and is switched on with
\code{\#FIELDSOLVER implicit1}. It differs from the recommendations above in two points.
\begin{itemize}
\item \textbf{Diagonal node volumes instead of $\Gm=\Pm^T\Wt\Pm$} (Q3). Following
  Lapenta (2017, eqs.~6 and 17), every unknown node has a volume
  $V_g=(\Pm^T\tilde{\mathbf w})_g$: one eighth of the volumes of the leaf cells around
  the node, plus the volumes of the hanging nodes that are interpolated from it,
  with the interpolation weights. The system is
  \begin{equation}
  \bigl[\mathsf I+\mathsf V^{-1}\bigl(\delta^2\Cm^T\Wb\Cm+4\pi\delta\,\Pm^T\widetilde\Mm\Pm\bigr)\bigr]\EE_\theta
  =\EE^n+\delta\,\mathsf V^{-1}\bigl(\Cm^T\Wb\BB^n-4\pi\Pm^T\hat{\tilde\jj}\bigr),
  \end{equation}
  and the conserved electric energy is $\sum_gV_g|\mathbf E_g|^2/8\pi$ over the unknown
  nodes. \code{calc\_E\_field\_energy} returns this sum when the composite solver is on.
  It differs from the corner average over the leaf cells by
  $\tilde w_h|\mathbf E_a-\mathbf E_b|^2/32\pi$ per hanging node $h$ on the edge $ab$.
\item \textbf{The coarser level owns the nodes on the interface} (Q1). The unknowns
  of level $\ell$ are the nodes that touch at least one leaf cell of level $\ell$ and,
  for $\ell>0$, are not level edge nodes. All the fine edge nodes, coincident and
  hanging, are then filled by one interpolation from the coarse level
  (\code{composite\_prolong}), and \code{composite\_fold} is its transpose. Both use
  the same weights on a copy of the coarse nodes with the fine box layout, so they are
  transposes by construction.
\end{itemize}
Everything else is as in Sections~\ref{sec:P}--\ref{sec:matvec}: own-level deposits
(\code{Pic::calc\_mass\_matrix} on every level, without the cross-level sums), the
curl--curl term as \code{curl\_center\_to\_node} of the masked
\code{curl\_node\_to\_center}, the right-hand side with the masked $\mathbf{B}$, and
the unchanged Faraday update. The covered coarse nodes are injected from the fine level
after the solve and the ghost nodes of the fine level are interpolated from the coarse
level, both for output only. GMRES uses the unweighted norm over the unknowns (P0 of
Section~\ref{sec:precond}). The solver needs
\code{ratioDivC2C}~$=0$, and stops with an error message with \code{\#SMOOTHE},
\code{\#SMOOTHJ}, \code{\#UPWINDE}, \code{\#COMOVING} or \code{\#BODY}. It also
stops with \code{\#UPWINDB} with \code{typeUpwindB curl}, with a time dependent
intrinsic magnetic field and with \code{\#EFIELDSOLVER} with \code{solverType amrex}.
With the developer switch \code{doTestComposite} of the \code{Pic} class (set in the
source), the symmetry of $\Pm^T\Cm_{\rm leaf}^T\Wb\Cm_{\rm leaf}\Pm$ is
checked with random vectors, and the energy budget of the first steps is printed.

\paragraph{Non-periodic boundaries (26 November 2026).} They are treated
exactly as in the level-by-level solver, by reusing its boundary routines
(this answers Q4 with the first option: a refined level may not touch a non-periodic
boundary, and the code stops if it does).
\begin{itemize}
\item The cells outside the domain count as leaf cells of the base level, so the
  boundary nodes keep the full stencil of $\Cm^T\Wb\Cm$ and the full volume $V_g$,
  as in the node based operator of the level solver.
\item The nodes outside the domain are filled by \code{apply\_field\_bc} when the
  operator acts on $\EE^n$, and by \code{apply\_field\_bc\_homogeneous} when it acts
  on the correction $\EE_\theta-\EE^n$ that GMRES solves for. The magnetic field in
  the cells outside the domain is filled as in \code{update\_E\_rhs}.
\item The rows of the nodes with a prescribed electric field (\code{inflow},
  \code{fixed}, \code{wave}) and of the tangential components on \code{conducting}
  walls are replaced by the identity, with the boundary value on the right-hand side.
\item The electric energy uses the node volumes of the cells inside the domain only,
  so that it agrees with the cell based sum of \code{calc\_E\_field\_energy}.
\end{itemize}
With these boundaries the total energy is not conserved: it changes by the flux
through the boundary. On a uniform grid, where the single solve can be forced for
testing, it reproduces the level solver to round-off for \code{outflow},
\code{conducting}, \code{absorb}, \code{fixed} and \code{inflow} boundaries. For
\code{wave} boundaries the two agree to round-off only if the level solver refreshes
the ghost nodes after \code{apply\_field\_bc}; as it is, the level solver uses the
old value of a boundary node in the ghost layer of the neighbouring box, a relative
difference of $10^{-6}$ in the test \code{bc\_wave}.

\paragraph{The other options of the level solver (26 November 2026).}
\code{\#FIELDSOLVER implicit1} now always uses the single solve, also on a uniform
grid, where it reproduces the level solver to round-off (to $10^{-9}$ with the
Newton--Krylov step of the nonlinear limiter of \code{\#UPWINDE}, the noise of its
finite difference). The options are added as follows.
\begin{itemize}
\item \code{ratioDivC2C} $=c>0$: the level solver uses
  $-\nabla^2+\nabla[(1-c)D_{\rm n2c}+c\,D_{\rm c2c}]$ for the curl curl. Here the
  curl curl is $\Cm^T\Wb\Cm$, so $c\,\nabla(D_{\rm c2c}-D_{\rm n2c})\EE$ of the leaf
  cells is added, weighted with the cell volume and folded like the other terms.
  The wider stencil uses the covered coarse nodes (injected from the fine level)
  and the ghost nodes of the fine levels (interpolated from the coarse level).
\item \code{\#UPWINDE}: the dissipation is added on every unknown node with the
  stencil of its own level, after the division by $V_g$.
\item \code{\#COMOVING} and the inductive field of a time dependent intrinsic
  magnetic field: $\pm V\,\Mm\EE_{\rm bg}$ is added to the right-hand side before
  the fold, like $\Mm\EE$ in the operator.
\item \code{\#SMOOTHJ}: the current is folded and divided by $V_g$ first, completed
  on the edge, covered and ghost nodes, and then smoothed level by level. The
  nodes inside a body keep the deposited values.
\item \code{\#SMOOTHE}: as in \code{update\_E\_impl}, the change of $\EE$ is
  smoothed, then $\EE_\theta$ and $\EE^{n+1}$ after the boundary conditions; the
  fine edge nodes are interpolated from the coarse nodes again afterwards.
\item \code{\#BODY}: the nodes of a line-tied body and the interior nodes of a
  conducting body are not unknowns; for a conducting body the rows and the
  right-hand side are projected on the radial direction on the surface.
\item \code{\#UPWINDB} with \code{typeUpwindB curl}: only the boundary filling of
  $\BB$ in the right-hand side differs; the update of $\BB$ is common.
\end{itemize}
None of the dissipative options conserves the energy, with either solver. Verified on
uniform grids against the level solver (20 steps, tolerance $10^{-11}$): all options
together on a periodic grid $10^{-15}$; outflow, conducting and fixed boundaries with
\code{ratioDivC2C} $=0.1$, \code{\#SMOOTHJ} and \code{\#UPWINDE} $10^{-15}$; line-tied,
conducting, insulating and reflecting bodies in 2D $10^{-14}$ with the linear limiter.
Two differences are not errors of the single solve: (i) \code{\#SMOOTHE} with
conducting walls and time dependent \code{wave} faces differ by $10^{-3}$ and
$10^{-7}$ unless the level solver refreshes the ghost nodes after
\code{apply\_field\_bc} in \code{update\_E\_matvec}; with that one line added to the
level solver the agreement is $10^{-15}$. (ii) For an intrinsic field that is not
periodic in a periodic domain the electric energy differs by $10^{-5}$, because the
node based sum takes the inductive field on one of the two periodic images of a face
node and the cell based sum on both; the solves are identical.

A refined level that touches a non-periodic boundary is still not allowed: the
boundary conditions are applied on all levels in the code, but the case could not be
tested, because the initial magnetic field is not a number there with any field
solver (also in the code before these changes, true 2D library, conducting or outflow
boundary). \code{\#EFIELDSOLVER} with \code{solverType amrex} and GPUs are not
supported either.

Cost on a uniform grid ($48^3$ cells, 8 particles per cell, 4 processes, 20 steps,
tolerance $10^{-11}$): the number of matrix-vector products is the same as with
\code{implicit}; the electric field solve takes 2\% more time without and 5\% more
with all the options above (1\% and 3\% of the run time).

The variant \code{implicit1} of the test \code{tests/amr\_equilibrium} checks that the
total energy is conserved to a relative $10^{-6}$.

%=====================================================================
\section{Open questions (summary)}\label{sec:open}
\begin{description}
\item[Q1] Coincident interface nodes owned by the fine (recommended) or the coarse level.
\item[Q2] Restriction $\Rm$ of the covered coarse $\mathbf{E}$: injection or full weighting (diagnostic only).
\item[Q3] Consistent $\Gm=\Pm^T\Wt\Pm$ (recommended; it is what the diagnostic reports) or lumped.
\item[Q4] Field BCs on refined levels that touch non-periodic boundaries.
\item[Q5] Preconditioner (P0$\to$P1/P2), GMRES norm (volume weighted) and tolerance.
\item[Q6] Nodal $\mathbf{B}$ at the interface for the pusher: composite average or ghost interpolation.
\item[Q7] 3D $\nabla\cdot\mathbf{B}$ at the interface: accept and clean, or a local
  composite divergence, or a modified $\Cm$.
\item[Q8] Gauss's law at the interface and a composite div~$\mathbf{E}$ correction.
\item[Q9] Proper nesting between levels $\ell+1$ and $\ell-1$.
\item[Q10] Features outside the identity even on one level (\code{coefDiff}, upwind
  $\mathbf{E}$, smoothing, co-moving frame, div~$\mathbf{E}$ correction): keep them excluded,
  or make each energy consistent (e.g.\ smoothing as $\Sm\to\Sm\mathsf F$ in gather,
  deposit and $\Mm$ together).
\end{description}

\end{document}
